\section*{Chapter \ref{ch:g12:catalysis} --- Catalysis}

\begin{solution}{exo:g12:catalysis:1}
\omterm{def:g7:chemical-reactions:reactant}{Reactant}: \ce{H2O2}. Products: \ce{H2O} and \ce{O2}. \omterm{def:g12:catalysis:catalyst}{Catalyst}: \ce{I-}
(used in step 1, given back in step 2). Intermediate: \ce{IO-} (formed in
step 1, used in step 2).
\end{solution}

\begin{solution}{exo:g12:catalysis:2}
Homogeneous; heterogeneous; heterogeneous; homogeneous.
\end{solution}

\begin{solution}{exo:g12:catalysis:3}
The uncatalysed path. The start and end levels are the same: the same
products form, with the same energy released; only the path differs.
\end{solution}

\begin{solution}{exo:g12:catalysis:4}
It is used in one step and given back in another: it appears on both
sides of the sum of the steps, and cancels.
\end{solution}

\begin{solution}{exo:g12:catalysis:5}
The reaction happens on the surface: the larger the surface, the more
\omterm{def:g7:atoms-and-molecules:molecule}{molecules} can react at once, for the same mass of (often expensive)
\omterm{def:g12:catalysis:catalyst}{catalyst}.
\end{solution}

\begin{solution}{exo:g12:catalysis:6}
Sum: \ce{2Fe^{3+} + 2H2O2 + 2Fe^{2+} + 2H+ -> 2Fe^{2+} + O2 + 2H+ +
2Fe^{3+} + 2H2O}. The iron \omterm{def:g9:ions:ion}{ions} and the hydrogen \omterm{def:g9:ions:ion}{ions} appear on both
sides and cancel: \ce{2H2O2 -> 2H2O + O2}.
\end{solution}

\begin{solution}{exo:g12:catalysis:7}
Half of \qty{72}{mL} is \qty{36}{mL}: reached after about \qty{35}{min}
without \omterm{def:g12:catalysis:catalyst}{catalyst} and \qty{3.5}{min} with it. The \omterm{def:g12:catalysis:catalyst}{catalyst} shortens it ten
times.
\end{solution}

\begin{solution}{exo:g12:catalysis:8}
The reaction takes place on the platinum surface; the lead covers it and
stays, so the gases can no longer reach the \omterm{def:g9:periodic-table-first-look:metal}{metal}. The \omterm{def:g12:catalysis:catalyst}{catalyst} is
``poisoned'': present but useless.
\end{solution}

\begin{solution}{exo:g12:catalysis:9}
Up to about \qty{37}{\celsius}, the reaction speeds up with the
temperature, like any reaction. Above, the \omterm{def:g12:catalysis:enzyme}{enzyme}, a protein, loses its
shape and with it its activity; at \qty{80}{\celsius} it is destroyed.
\end{solution}

\begin{solution}{exo:g12:catalysis:10}
\ce{2CO + 2NO -> 2CO2 + N2}: C 2 = 2, O 4 = 4, N 2 = 2. Carbon monoxide is
oxidised by nitrogen monoxide, which is reduced: each pollutant removes
the other.
\end{solution}

\begin{solution}{exo:g12:catalysis:11}
$n(\ce{O2}) = 0.072 / 24.1 = \qty{2.99e-3}{mol}$;
$n(\ce{H2O2}) = 2 \times \num{2.99e-3} = \qty{5.98e-3}{mol}$ in
\qty{10.0}{mL}: $c = \qty{0.598}{mol/L}$.
\end{solution}

\begin{solution}{exo:g12:catalysis:12}
\ce{C2H5OH -> C2H4 + H2O}, an \omterm{def:g12:curly-arrows:categories}{elimination}; \ce{C2H5OH -> CH3CHO + H2}.
A \omterm{def:g12:catalysis:selective}{selective catalyst} speeds up one of several possible reactions much
more than the others, and so decides the product.
\end{solution}

\begin{solution}{exo:g12:catalysis:13}
The \omterm{def:g12:catalysis:catalyst}{catalyst} does not change the \omterm{def:g10:reaction-progress-table:system}{final state}: the curves with and
without \omterm{def:g12:catalysis:catalyst}{catalyst} end at the same volume of oxygen. The same amount of
product is obtained, only sooner.
\end{solution}

\begin{solution}{exo:g12:catalysis:14}
$\num{6.02e23} / \num{5e6} = \qty{1.2e17}{s}$, some four billion years.
To do it in one second: \num{1.2e17} \omterm{def:g12:catalysis:enzyme}{enzyme} \omterm{def:g7:atoms-and-molecules:molecule}{molecules}, that is
\qty{2e-7}{mol}.
\end{solution}

\begin{solution}{exo:g12:catalysis:15}
Ammonia \ce{NH3}: $150 \times 17.0 / 14.0 = \qty{182}{Mt}$. Nitrogen
\omterm{def:g7:atoms-and-molecules:atom}{atoms}: $\num{1.50e14} / 14.0 = \qty{1.07e13}{mol}$, that is
\qty{5.36e12}{mol} of \ce{N2}.
\end{solution}

\begin{solution}{pb:g12:catalysis:1}
\textbf{1.} \ce{2H2O2 -> 2H2O + O2}.

\textbf{2.} Yes: hydrogen peroxide is the \omterm{def:g11:redox:oxidant}{oxidant} of the couple
\ce{H2O2}/\ce{H2O} and the \omterm{def:g11:redox:oxidant}{reductant} of the couple \ce{O2}/\ce{H2O2}; it
oxidises and reduces itself.

\textbf{3.} Without a \omterm{def:g12:catalysis:catalyst}{catalyst} the reaction is extremely slow.

\textbf{4.} $20 / 22.4 = \qty{0.893}{mol}$.

\textbf{5.} Two \omterm{def:g10:the-mole:amount}{moles} of hydrogen peroxide per \omterm{def:g10:the-mole:amount}{mole} of oxygen:
\qty{1.79}{mol}.

\textbf{6.} \qty{1.79}{mol/L}.

\textbf{7.} $M = 2 \times 1.0 + 2 \times 16.0 = \qty{34.0}{g/mol}$;
$C_m = 1.79 \times 34.0 = \qty{60.7}{g/L}$.

\textbf{8.} Half: \qty{0.893}{mol/L}.

\textbf{9.} Catalase: an \omterm{def:g12:catalysis:enzyme}{enzyme}, dissolved, a biological \omterm{def:g12:catalysis:catalyst}{catalyst};
iron(III) \omterm{def:g9:ions:ion}{ions}: \omterm{def:g12:catalysis:homogeneous}{homogeneous catalysis}.

\textbf{10.} \ce{2Fe^{3+} + H2O2 -> 2Fe^{2+} + O2 + 2H+} and
\ce{2Fe^{2+} + H2O2 + 2H+ -> 2Fe^{3+} + 2H2O}; their sum is
\ce{2H2O2 -> 2H2O + O2}.

\textbf{11.} With a \omterm{def:g12:catalysis:catalyst}{catalyst} the volume rises much faster; both curves
level off at the same final volume.

\textbf{12.} The iron(III) \omterm{def:g9:ions:ion}{ions} (orange) are turned into iron(II) (pale
green) in the first step and back into iron(III) in the second; at the
end, all the iron is iron(III) again.

\textbf{13.} Boiling has destroyed the \omterm{def:g12:catalysis:enzyme}{enzyme}.

\textbf{14.} $0.100 \times 20 = \qty{2.0}{L}$.

\textbf{15.} $n(\ce{O2}) = 0.100 \times 0.893 = \qty{0.0893}{mol}$, that
is $0.0893 \times 24.1 = \qty{2.15}{L}$.

\textbf{16.} The slow decomposition releases oxygen, which would build up
pressure in a tight bottle.

\textbf{17.} $1.50 / 2 = \qty{0.750}{mol}$ of oxygen per litre, that is
$0.750 \times 22.4 = \qty{16.8}{L}$: ``16.8 volumes''.

\textbf{18.} $(1.79 - 1.50)/1.79 \approx \qty{16}{\percent}$.

\textbf{19.} \qty{1.79}{mol/L}.
\end{solution}
